% 调用上级目录公共模板
\documentclass{../template/softeng}

\begin{document}

% 封面：文档名、项目名、版本
\doccover{软件需求规格说明书}
{校园图书借阅管理系统}
{V2026.07}

% 目录
\tableofcontents
\newpage

% 修订记录
\begin{revrecord}
V1.0 & 2026-07-13 & 编写人 & 审核人 & 初稿完成，完成功能性需求与非功能性需求 \\
\hline
\end{revrecord}

% ============================================================
%  1. 引言
% ============================================================
\section{引言}
\subsection{编写目的}
本文档为校园图书借阅管理系统软件需求规格说明书，用于完整定义系统全部功能性与非功能性需求，明确项目开发边界、业务规则、数据规范、运行约束。
本文档交付对象包含项目经理、系统设计人员、前后端开发工程师、测试人员，作为项目规划、架构设计、代码开发、功能测试、验收交付的核心依据。

\subsection{项目背景}
本系统为高校大二软件工程实训项目，开发单位为在校实训开发小组，服务对象为高校图书馆管理员与在校师生。
当前高校图书馆仍存在借阅流程人工化、书目查询效率低、超期提醒不及时等问题。本系统独立部署运行，依托Web服务与关系型数据库实现完整借阅业务闭环，用于支撑图书在线检索、自助借还、超期管理、阅览统计等功能，提升图书馆管理效率与师生借阅体验。

\subsection{定义}
\begin{itemize}
    \item ISBN：International Standard Book Number，国际标准书号；
    \item 索书号：图书馆为图书分配的唯一排架编号；
    \item 借阅证：读者在系统中的唯一身份标识；
    \item 超期：借阅时间超过规定借阅期限的状态；
    \item 续借：在借阅到期前延长借阅期限的操作；
    \item 馆藏：图书馆实际拥有的图书副本数量。
\end{itemize}

\subsection{参考资料}
\begin{enumerate}
    \item GB/T 9385-2008 计算机软件需求规格说明规范；
    \item 实训项目任务书（2026软件工程实训）；
    \item Web前后端分离开发规范文档；
    \item 高校图书馆管理信息系统技术标准。
\end{enumerate}

% ============================================================
%  2. 任务概述
% ============================================================
\section{任务概述}
\subsection{目标}
搭建一套Web端校园图书借阅管理平台，区分读者、图书管理员、系统管理员三类角色；实现图书在线检索、自助借阅与归还、超期自动提醒、借阅历史查询、热门图书排行功能；管理员端提供图书入库编目、读者档案管理、借阅规则配置、超期罚款处理、流通数据统计功能，满足高校图书馆日常借阅管理与阅览服务需求。

\subsection{运行环境}
\begin{itemize}
    \item 服务操作系统：Windows Server / Linux CentOS；
    \item Web服务：Nginx + Spring Boot；
    \item 数据库：MySQL 8.0；
    \item 前端框架：Vue 3 + Element Plus；
    \item 客户端：主流Chrome、Edge浏览器。
\end{itemize}

\subsection{条件与限制}
\begin{enumerate}
    \item 系统需支持至少200名用户同时在线检索图书；
    \item 单次借阅操作响应时间不超过2秒；
    \item 读者同时借阅上限为5本，借阅期限为30天，可续借1次（15天）；
    \item 仅支持校内师生凭学号/工号注册，不开放校外读者；
    \item 系统需与学校统一认证系统对接（学号/工号单点登录）。
\end{enumerate}

% ============================================================
%  3. 数据描述
% ============================================================
\section{数据描述}
\subsection{静态数据}
系统长期持久存储不变基础数据：读者账号信息、图书元数据（书名、作者、ISBN、出版社、索书号）、书架位置信息、分类标签、借阅规则配置参数、图书馆基础信息。

\subsection{动态数据}
系统运行过程中实时产生、变更的数据：读者登录会话、借阅记录、归还记录、续借记录、超期罚款记录、图书预约队列、热门图书排行榜、流通统计数据、读者行为日志。

\subsection{数据库介绍}
采用MySQL 8.0关系型数据库存储全部业务数据，划分读者表（reader）、图书表（book）、借阅记录表（borrow\_record）、罚款记录表（fine）、预约记录表（reservation）、分类表（category）、系统日志表（sys\_log），数据表通过外键建立关联约束。

\subsection{数据词典}
包含读者ID（学号/工号）、姓名、角色类型、图书ISBN、书名、作者、出版社、索书号、馆藏数量、可借数量、借阅日期、应还日期、实际归还日期、续借次数、罚款金额、罚款状态等全部字段含义、数据类型、取值范围说明。

\subsection{数据采集}
\begin{itemize}
    \item 基础图书数据：管理员通过后台表单录入或批量Excel导入；
    \item 读者数据：对接学校统一认证系统，自动同步师生信息；
    \item 借阅数据：读者扫码/检索借阅时实时写入借阅记录；
    \item 统计数据：系统定时任务每日凌晨汇总流通量、热门图书、超期率。
\end{itemize}

% ============================================================
%  4. 功能需求
% ============================================================
\section{功能需求}
\subsection{功能划分}
系统整体划分为五大功能模块：读者自助服务模块、图书检索模块、借阅管理模块、后台管理模块、统计报表模块。

\subsection{功能描述}

% --- 读者自助服务 ---
\reqitem{FR001}{读者可通过学号/工号登录系统，查看个人借阅记录、当前借阅状态与到期提醒。}
\reqitem{FR002}{读者可在线自助续借未到期图书（每本限续借1次，延长15天）。}
\reqitem{FR003}{读者可预约已被借出的图书，图书归还后自动通知预约读者。}

% --- 图书检索 ---
\reqitem{FR004}{读者可按书名、作者、ISBN、索书号、分类进行图书检索，支持模糊搜索。}
\reqitem{FR005}{图书详情页展示书目信息、馆藏位置、当前可借数量、借阅热度。}
\reqitem{FR006}{系统展示热门借阅排行、新书上架推荐，辅助读者发现图书。}

% --- 借阅管理 ---
\reqitem{FR007}{图书管理员扫描读者借阅证与图书条码，完成借阅登记，系统自动计算应还日期。}
\reqitem{FR008}{读者归还图书时，管理员扫描图书条码完成归还，系统自动检测是否超期并计算罚款。}
\reqitem{FR009}{系统自动检测超期未还图书，向读者发送催还通知（站内消息+邮件）。}

% --- 后台管理 ---
\reqitem{FR010}{图书管理员可新增、编辑、下架图书信息，支持批量Excel导入馆藏书目。}
\reqitem{FR011}{系统管理员可管理读者账号、配置借阅规则（借阅上限、期限、续借次数、罚款标准）。}
\reqitem{FR012}{系统管理员可查看操作日志，追溯所有借还操作与后台管理行为。}

% --- 统计报表 ---
\reqitem{FR013}{系统自动生成图书流通量日报/月报、分类借阅占比、超期率统计。}
\reqitem{FR014}{系统生成读者借阅排行榜、热门图书Top20，支撑图书馆采购决策。}

% ============================================================
%  5. 性能需求
% ============================================================
\section{性能需求}
\subsection{数据精确度}
借阅日期精确至秒，罚款金额精确至小数点后两位；图书检索结果按相关度排序，无漏检、错检；统计数据与原始记录严格一致，无统计偏差。

\subsection{时间特性}
\begin{itemize}
    \item 网页页面加载响应时间 $\le$ 1.5s；
    \item 图书检索查询响应时间 $\le$ 1s；
    \item 借阅/归还操作处理时间 $\le$ 2s；
    \item 批量导入1000条图书数据耗时 $\le$ 30s；
    \item 统计报表生成时间 $\le$ 5s。
\end{itemize}

\subsection{适应性}
系统可适配Windows、Linux服务端环境；支持切换MySQL/PostgreSQL数据库；可扩展支持RFID自助借还设备接口；可调整借阅规则参数配置（期限、上限、罚款标准）。

% ============================================================
%  6. 运行需求
% ============================================================
\section{运行需求}
\subsection{用户界面}
采用Web网页可视化界面，响应式布局适配PC端与移动端。读者端：图书检索页（分类筛选+搜索框）、图书详情页（封面+书目信息+可借状态）、个人中心（借阅列表+到期提醒）。管理端：后台仪表盘（借阅总量+超期数+实时数据）、书目管理表格（增删改查+批量导入）、读者管理列表、流通数据统计图表。

\subsection{硬件接口}
\begin{itemize}
    \item 条码扫描枪：USB接口，用于扫描读者借阅证条码与图书ISBN条码；
    \item 打印机：USB/网络接口，用于打印罚款收据与催还通知单。
\end{itemize}

\subsection{软件接口}
\begin{itemize}
    \item 数据库接口：MySQL标准JDBC连接接口；
    \item 前端接口：前后端RESTful API数据交互接口，JSON格式；
    \item 学校统一认证接口：CAS/OAuth 2.0协议，实现学号/工号单点登录；
    \item 邮件接口：SMTP协议，用于发送催还通知与预约到馆通知。
\end{itemize}

\subsection{故障处理}
\begin{itemize}
    \item 数据库断开：前端页面提示"服务暂不可用"，缓存当前操作上下文；
    \item 网络超时：借阅操作采用事务回滚机制，确保数据一致性；
    \item 条码扫描失败：弹窗提示手动输入借阅证号或ISBN；
    \item 并发冲突（同本书多人预约）：按预约时间先后排队，先到先得。
\end{itemize}

% ============================================================
%  7. 其它需求
% ============================================================
\section{其它需求}
\begin{enumerate}
    \item 安全性：区分读者、图书管理员、系统管理员三级角色权限；密码采用BCrypt加密存储；SQL参数化查询防止注入攻击；敏感操作记录操作日志；
    \item 可维护性：后台提供可视化管理页面，代码遵循MVC分层架构，关键模块有单元测试覆盖；
    \item 可移植性：基于Java Spring Boot + Vue前后端分离架构，可快速部署至不同服务器环境；
    \item 易用性：读者检索借阅流程不超过3步操作，管理后台提供批量操作入口减少重复劳动；
    \item 可扩展性：预留RFID接口、微信小程序接口、图书推荐算法模块扩展点。
\end{enumerate}

% ============================================================
%  附录：用例图
% ============================================================
\begin{figure}[htbp]
    \centering
    \includegraphics[width=0.7\textwidth]{usecase.png}
    \caption{校园图书借阅管理系统整体用例图（待替换为实际用例图）}
\end{figure}

\end{document}
